\documentclass[11pt]{article}

\usepackage[T1]{fontenc}
\usepackage{lmodern}
\usepackage[a4paper,margin=29mm]{geometry}
\usepackage{amsmath,amssymb,amsthm,mathtools}
\usepackage{microtype}
\usepackage{enumitem}
\usepackage{graphicx}
\usepackage{xcolor}
\usepackage[colorlinks=true,linkcolor=blue!45!black,citecolor=blue!45!black,
  urlcolor=blue!55!black]{hyperref}
\setlength{\emergencystretch}{2em}

\newtheorem{theorem}{Theorem}[section]
\newtheorem{proposition}[theorem]{Proposition}
\newtheorem{lemma}[theorem]{Lemma}
\newtheorem{corollary}[theorem]{Corollary}
\newtheorem{definition}[theorem]{Definition}
\newtheorem{remark}[theorem]{Remark}

\newcommand{\C}{\mathbb C}
\newcommand{\R}{\mathbb R}
\newcommand{\N}{\mathbb N}
\newcommand{\cZ}{\mathcal Z}
\newcommand{\cG}{\mathcal G}
\newcommand{\cP}{\mathcal P}
\newcommand{\cA}{\mathcal A}
\newcommand{\cE}{\mathcal E}
\newcommand{\Res}{\operatorname*{Res}}
\newcommand{\ord}{\operatorname{ord}}
\newcommand{\supp}{\operatorname{supp}}

\title{A Local-to-Global Dichotomy for the Riemann Hypothesis:\\
Universal Finite-Factor Orbits and One-Sequence Xi-Divisor Tomography}
\author{Michael Haufschild}
\date{}

\begin{document}

\maketitle

\begin{abstract}
We prove an unconditional local-to-global structure theorem for source approaches to the
Riemann hypothesis.  Its local half is spectrum-independent.  Let
\(\Lambda=(\lambda_j)\) be an injective positive frequency family such that, for some
\(R_0>0\), every integer \(m\ge1\), every \(R\ge R_0\), and after every finite deletion,
\[
 \sum_{\lambda_j>R+1}\frac1{(\lambda_j-R)^m+1}=\infty.
\]
On every closed disc centred at \(1/2\) inside the critical strip, the pole-free additive
orbit of arbitrarily small reflected local-factor logarithmic derivatives is dense in the
full holomorphic, functional-equation-odd, Schwarz-real fibre.  Its maximal continuous
Hausdorff quotient is therefore one point.  Prime logarithms satisfy the hypothesis by
Euler's divergence of \(\sum_p p^{-1}\).  Consequently no continuous compact-open
local-germ reader which ignores finite Euler-factor changes can classify an RH-equivalent
Hankel sign.

The global half constructs the missing survivor.  Fix any width \(L>0\), translate one
symmetric two-edge window of that width, and sample its Xi--Weil polarization at
\(a=\frac12\log n\), obtaining \(q_L(n)\).  The ordinary Dirichlet series of this single
sequence continues meromorphically to \(\Re s>1\).  Its poles are exactly \(1+w\), for the
off-critical centred Xi zeros with \(\Re w>0\), and normalized residues recover their
analytic multiplicities.  Thus one sequence reconstructs the complete off-critical divisor.
RH is equivalent to holomorphy in that half-plane, to nonnegativity of \(q_L\), to eventual
lower boundedness, and to boundedness.  If RH is false and
\[
 \delta_*=\sup_{\Xi(\rho)=0,\ \Re\rho>1/2}(\Re\rho-\tfrac12),
\]
then \(n^{-\alpha}q_L(n)\) is cofinally unbounded in both signs for every
\(0\le\alpha<\delta_*\).

Every admissible sample has the exact source evaluation
\[
 q_L(n)=g_L(n)
 =-2\sum_{2\le k\le n}\frac{\Lambda(k)}{\sqrt{k}}\,
    R_{L,\frac12\log n}(\log k)
   +\Re\cA_{\infty,L,\frac12\log n},
\]
with a finite von Mangoldt sum and an absolutely convergent Gamma--polar integral.
After any admissible finite-prefix deletion, the negative part has exact Dirichlet abscissa
\[
 \sigma_{\rm abs}(g^-_{L,N})=
 \begin{cases}
 -\infty,&\mathrm{RH},\\
 1+\delta_*,&\neg\mathrm{RH}.
 \end{cases}
\]
Every horizontal zero-free region is thereby equivalent to a cumulative negative-mass
growth law.  Individual completeness also holds for every profile in a dense residual
smooth class.  No rightmost zero, simplicity, zero gap, numerical zero information, or
finiteness hypothesis on off-line zeros is used.

This does not prove RH.  It proves a sharp dichotomy: compact-open local continuation
annihilates continuous finite-factor-insensitive RH information, whereas one explicit
global weighted-tail sequence retains the entire off-critical divisor and quantitatively
classifies every possible failure.
\end{abstract}

\medskip
\noindent\textbf{2020 Mathematics Subject Classification.}
Primary 11M26; Secondary 11M06, 30B50.

\noindent\textbf{Keywords.}
Riemann hypothesis; Weil criterion; explicit formula; completeness; Euler factors;
Dirichlet series.

\section{Introduction}

Let
\[
 \Xi(s)=\frac12s(s-1)\pi^{-s/2}\Gamma(s/2)\zeta(s).
 \tag{1.1}
\]
Weil's positivity criterion characterizes RH by nonnegativity of a quadratic functional on a
large test-function class \cite{Weil1952,Bombieri2000}.  Later work has studied localized Weil
forms, screw functions, semilocal trace formulae, and finite spectral models
\cite{Connes1999,ConnesConsani2021,Suzuki2023,Suzuki2025,Suzuki2026,CCM2025,Groskin2026}.
Those formulations raise a basic compression problem: how much of the test-function class is
needed to retain the horizontal Xi divisor?

There is a logically prior information question.  The prime and archimedean source can be
continued to a holomorphic germ at the functional-equation centre.  Can a continuous reader of
that local germ, insensitive to finitely many Euler factors, retain the sign information needed
for RH?  We prove that the answer is no for a spectrum-independent reason.  Reflected
derivative-Lambert local factors generate a dense orbit in the entire analytic odd germ fibre;
for prime logarithms, this remains true after arbitrary finite prime deletion.  Hence the
corresponding continuous invariant quotient is trivial.

The fixed-width sequence in this paper supplies the exact positive counterpart.  It is not a
function of the compact-open centre germ alone: it retains the logarithmically indexed,
place-weighted source tail.  That extra global presentation is sufficient to reconstruct every
off-critical zero and its multiplicity.  The paper therefore proves both sides of a
local-to-global dichotomy rather than merely recommending that a successful method be
``nonlocal.''

Yoshida's and Suzuki's localization theorems test an entire function space while the interval
expands.  Bombieri related negative eigenspaces of localized forms to off-line zeros, with a
finiteness hypothesis in the limiting theorem and a possible linear-dependence alternative
\cite{Yoshida1992,Bombieri2000,Bombieri2003,Suzuki2026}.  The construction below uses neither an
optimized vector nor an eigenfunction.  One rectangle of any prescribed positive width is fixed
in advance; only its translation changes.  If RH is false, that single scalar observable has
arbitrarily late negative and positive excursions at every power below the horizontal support
endpoint.  Thus no finite-off-line-zero hypothesis, density premise, or limiting
linear-independence condition remains.

This paper gives an exact answer for one elementary geometry.  Width is fixed once and for all.
Only aperture moves, and only on the logarithmic integer mesh.  The resulting scalar sequence
does not merely decide the critical-line case.  Its Dirichlet continuation recovers every
off-critical zero and its multiplicity.  Its lower-loss exponent, negative Dirichlet abscissa,
and cumulative negative-mass exponent recover every symmetric horizontal zero-free strip.

The global result is a classification theorem, not a new assertion of RH.  Its logical content is
\[
\begin{gathered}
 \text{off-critical Xi divisor with multiplicity}\\
 \Updownarrow\\
 \text{polar divisor of one fixed-width Weil sequence}\\
 \Updownarrow\\
 \text{one explicit finite prime--Gamma sequence},
\end{gathered}
\tag{1.2}
\]
with its lower-loss and negative-mass spectra furnishing the corresponding horizontal support
law.
False RH has a compulsory signature: at every power below the transverse support endpoint, the
same observable makes arbitrarily late excursions of both signs.  This rules out a common
weakness of RH equivalents, namely that failure is encoded only as the negation of an unquantified
positivity statement.

Together with the local orbit theorem, this gives the information diagram
\[
\begin{array}{ccc}
\begin{gathered}\text{continued centre germ}\\[-1mm]\text{(compact-open)}\end{gathered}
 &\longrightarrow&
\begin{gathered}\text{one-point quotient}\\[-1mm]\text{(finite-factor invariant)}\end{gathered}\\[2mm]
\rotatebox{90}{\(\not\!\Longleftarrow\)}
&&\\[-1mm]
\begin{gathered}\text{weighted logarithmic}\\[-1mm]\text{prime--Gamma tail}\end{gathered}
 &\longrightarrow&
\text{off-critical Xi divisor}.
\end{array}
\tag{1.2a}
\]
The lower arrow is realized by a single explicit scalar sequence.  Thus the theorem both
identifies what local continuation loses and constructs the minimal kind of global datum that
survives.

\subsection{Main definitions}

Write each nontrivial zero as \(\rho=\frac12+w\), retaining analytic multiplicity, and let
\(\cZ\) be the centered multiset of all such \(w\).  For \(L>0\), put
\[
 A_L(w)=
 \begin{cases}
  (1-e^{-Lw})/w,&w\ne0,\\
  L,&w=0,
 \end{cases}
 \qquad
 F_{a,L}(w)=e^{aw}A_L(w)+e^{-aw}A_L(-w).
 \tag{1.3}
\]
This is the bilateral Laplace transform of
\[
 h_{L,a}(x)=\mathbf 1_{[a-L,a]}(x)+\mathbf 1_{[-a,-a+L]}(x).
 \tag{1.4}
\]
The fixed-width Weil observable is the absolutely convergent divisor polarization
\[
 Q_L(a)=\sum_{w\in\cZ}F_{a,L}(w)^2,
 \qquad
 q_L(n)=Q_L\!\left(\tfrac12\log n\right)\quad(n\ge1).
 \tag{1.5}
\]
The functional equation and Schwarz symmetry make \(Q_L(a)\) real.  An integer \(n\) is
\emph{\(L\)-admissible} when \(L\le\frac12\log n\), and \emph{strictly \(L\)-admissible} when
\(L<\frac12\log n\).

Let
\[
 \delta_*=\sup\{\Re w:w\in\cZ,\ \Re w>0\}
 \tag{1.6}
\]
when an off-line zero exists, and set \(\delta_*=0\) under RH.  Thus
\(0<\delta_*\le\frac12\) under false RH.

\subsection{Main theorem}

\begin{theorem}[One-sequence divisor tomography]\label{thm:divisor-main}
Fix \(L>0\), and initially set
\[
 \mathcal D_L(s)=\sum_{n\ge1}\frac{q_L(n)}{n^s}
 \qquad(\Re s>1+\delta_*).
\]
Then \(\mathcal D_L\) has a canonical meromorphic continuation
\(\mathcal D_L^\#\) to \(\Re s>1\).  In that half-plane:
\begin{enumerate}[label=\textup{(\roman*)}]
\item its poles are exactly \(s=1+w\), where \(w\in\cZ\) and \(\Re w>0\);
\item every such pole is simple, and if \(w\) has analytic multiplicity \(m(w)\), then
\[
 \Res_{s=1+w}\mathcal D_L^\#(s)=2m(w)A_L(w)^2;
\]
\item normalized residue recovers multiplicity:
\[
 m(w)=\frac{\Res_{s=1+w}\mathcal D_L^\#(s)}{2A_L(w)^2}.
\]
\end{enumerate}
The functional equation and Schwarz symmetry recover the complete off-critical Xi divisor from
this right-half divisor.  In particular,
\[
 \mathrm{RH}
 \Longleftrightarrow
 \mathcal D_L^\#\text{ is holomorphic throughout }\Re s>1.
\]
\end{theorem}

\begin{theorem}[Fixed-width Xi tomography]\label{thm:main}
Fix \(L>0\).
\begin{enumerate}[label=\textup{(\roman*)}]
\item The following statements are equivalent:
\[
\begin{split}
 &\mathrm{RH};\\
 &q_L(n)\ge0\quad\text{for every positive integer }n;\\
 &q_L(n)\text{ is bounded below on an admissible tail};\\
 &|q_L(n)|\text{ is bounded on the admissible mesh}.
\end{split}
\tag{1.7}
\]
\item If RH is false, then for every \(0\le\alpha<\delta_*\), \(C\ge0\), and \(N_0\), there
exist \(m,n\ge N_0\), both \(L\)-admissible, such that
\[
 q_L(m)<-C m^\alpha,
 \qquad
 q_L(n)>C n^\alpha.
 \tag{1.8}
\]
Moreover \(|q_L(n)|=O_L(n^{\delta_*})\).  Hence \(\delta_*\) is the exact two-sided power
threshold.
\item For every \(\alpha\ge0\),
\[
\begin{split}
 &\exists C,N_0\ \forall L\text{-admissible }n\ge N_0,\quad
 q_L(n)\ge-Cn^\alpha\\
 &\hspace{28mm}\Longleftrightarrow\quad
 \Xi(\rho)=0\Longrightarrow \Re\rho\le\tfrac12+\alpha.
\end{split}
\tag{1.9}
\]
\end{enumerate}
\end{theorem}

Part (i) compresses Weil positivity to one fixed width and one discrete translation parameter.
Parts (ii)--(iii) are stronger than a binary equivalent: they classify the magnitude and sign of
every possible off-line defect.  The profile is explicit and fixed before any information about
the zeros is supplied; no minimization over a localized Hilbert space occurs.

\begin{theorem}[Exact finite prime--Gamma spectrum]\label{thm:source-main}
Fix \(L>0\) and a cutoff \(N\) for which every \(n\ge N\) is \(L\)-admissible.  There is an
explicit source coefficient \(g_L(n)\), consisting at each \(n\) of a finite von Mangoldt sum and
one absolutely convergent Gamma--polar integral, such that
\[
 g_L(n)=q_L(n)\qquad(n\ge N).
\tag{1.10}
\]
Put \(g^-_{L,N}(n)=\max\{-g_L(n),0\}\) for \(n\ge N\), and \(0\) otherwise.  Then
\[
 \sigma_{\rm abs}(g^-_{L,N})=
 \begin{cases}
  -\infty,&\mathrm{RH},\\
  1+\delta_*,&\neg\mathrm{RH}.
 \end{cases}
\tag{1.11}
\]
For every \(\alpha\ge0\), the following are equivalent:
\begin{enumerate}[label=\textup{(\alph*)}]
\item every Xi zero satisfies \(\Re\rho\le\frac12+\alpha\);
\item for every \(\varepsilon>0\),
\[
 \sum_{n\ge1}\frac{g^-_{L,N}(n)}
 {n^{1+\alpha+\varepsilon}}<\infty;
 \tag{1.12}
\]
\item for every \(\varepsilon>0\),
\[
 \sum_{1\le n\le X}g^-_{L,N}(n)
 =O_{\varepsilon,L,N}(X^{1+\alpha+\varepsilon}).
 \tag{1.13}
\]
\end{enumerate}
Under false RH, (1.13) holds with \(\alpha=\delta_*\), while for every
\(0<\eta<1+\delta_*\),
\[
 \sum_{1\le n\le X}g^-_{L,N}(n)
 \ne O(X^{1+\delta_*-\eta}).
\tag{1.14}
\]
In particular,
\[
 \mathrm{RH}\Longleftrightarrow
 \sum_{1\le n\le X}g^-_{L,N}(n)
 =O_{\varepsilon,L,N}(X^{1+\varepsilon})
 \quad\text{for every }\varepsilon>0.
\tag{1.15}
\]
\end{theorem}

The last equivalence does not supply an RH proof: its critical source estimate is precisely the
endpoint selected by the theorem.  The unconditional advance is the complete classification
(1.11)--(1.14), including the exact exponent forced by any failure.

\begin{corollary}[Prime--Gamma divisor tomography]\label{cor:source-divisor}
With \(L\) and \(N\) as in Theorem \ref{thm:source-main}, initially put
\[
 \mathcal G_{L,N}(s)=\sum_{n\ge N}\frac{g_L(n)}{n^s}.
\]
This series has a meromorphic continuation to \(\Re s>1\) whose poles and residues are exactly
those of \(\mathcal D_L^\#\).  Consequently the one explicit finite prime--Gamma coefficient
sequence determines every off-critical Xi zero and its analytic multiplicity, and RH is
equivalent to holomorphy of its continuation throughout \(\Re s>1\).
\end{corollary}

\begin{theorem}[Residual smooth robustness]\label{thm:smooth-main}
Let \(0<\eta<L-\eta\).  In the Fr\'echet space of smooth real profiles supported in
\([\eta,L-\eta]\), there is a dense residual subset \(\cG_{\eta,L}\) with the following property.
If RH is false, then for every \(f\in\cG_{\eta,L}\), \(B\in\R\), and \(N_0\), some
strictly \(L\)-admissible \(n\ge N_0\) satisfies
\[
 Q_{f,L}\!\left(\tfrac12\log n\right)<B.
\tag{1.16}
\]
For every single \(f\in\cG_{\eta,L}\), the following are equivalent:
\[
\begin{aligned}
\mathrm{RH}
&\Longleftrightarrow
Q_{f,L}\!\left(\tfrac12\log n\right)\ge0
\quad\text{at every strictly admissible sample}\\
&\Longleftrightarrow
\text{that one sequence is eventually bounded below}.
\end{aligned}
\tag{1.17}
\]
Consequently the corresponding class-wide nonnegativity and lower-boundedness criteria are also
equivalent to RH.
\end{theorem}

\begin{theorem}[Universal finite-factor orbit]\label{thm:local-orbit}
Let \(\Lambda=(\lambda_j)_{j\in J}\subset(0,\infty)\) be injective and suppose that, for
some \(R_0>0\), every integer \(m\ge1\), and every \(R\ge R_0\),
\[
 \sum_{\lambda_j>R+1}\frac1{(\lambda_j-R)^m+1}=\infty.
\tag{1.18}
\]
For \(0\le r<r_1<1/2\), put
\[
 K_r=\{s:|s-\tfrac12|\le r\},\qquad
 U_\lambda(s)=
 \lambda\frac{e^{-\lambda s}}{(1-e^{-\lambda s})^2}
 -
 \lambda\frac{e^{-\lambda(1-s)}}{(1-e^{-\lambda(1-s)})^2}.
\tag{1.19}
\]
Then the real span of the \(U_{\lambda_j}\) is dense, uniformly on \(K_r\), in every
holomorphic germ \(f\) satisfying
\[
 f(1-s)=-f(s),\qquad f(\overline s)=\overline{f(s)}.
\tag{1.20}
\]
For \(b\in\R\), let
\[
\begin{aligned}
 \Phi_{\lambda,a}(s)&=\lambda\frac{ae^{-\lambda s}}{1-ae^{-\lambda s}},\\
 \Delta_{\lambda,b}(s)
 &=\bigl[\Phi_{\lambda,1+b}(s)-\Phi_{\lambda,1}(s)\bigr]
  -\bigl[\Phi_{\lambda,1+b}(1-s)-\Phi_{\lambda,1}(1-s)\bigr].
\end{aligned}
\tag{1.21}
\]
For every \(\epsilon>0\), the additive group generated by those
\(\Delta_{\lambda_j,b}\) with
\[
 |b|<\epsilon,\qquad
 \sup_{s\in K_r}|(1+b)e^{-\lambda_j s}|<1
\tag{1.22}
\]
has the entire analytic odd Schwarz-real fibre in its uniform closure.  Consequently every
continuous Hausdorff-valued invariant of this pole-free group is constant on that fibre.
All conclusions survive arbitrary finite deletion from \(\Lambda\).

For \(\lambda_p=\log p\), hypothesis (1.18) holds after every finite prime deletion.
\end{theorem}

\begin{theorem}[Local-to-global RH dichotomy]\label{thm:dichotomy}
Let
\[
 \Xi_C(z)=\Xi(\tfrac12+iz),\qquad
 G_\zeta(z)=-\frac1z\frac{\Xi_C'(z)}{\Xi_C(z)}
\tag{1.23}
\]
at the removable point \(z=0\).  If
\(G_\zeta(z)=\sum_{n\ge0}M_nz^{2n}\), then
\[
 \mathrm{RH}\quad\Longleftrightarrow\quad
 \sum_{i,j}c_ic_jM_{i+j}\ge0
 \quad\text{for every }(c_0,\ldots,c_d)\in\mathbb Z^{d+1}
\tag{1.24}
\]
and every \(d\ge0\).

After any finite prime deletion, there is no continuous compact-open reader of the critical
centre germ which is invariant under the pole-free local-factor group of
Theorem~\ref{thm:local-orbit} and classifies the positivity in (1.24) on all analytic even
Schwarz-real moment germs.  In contrast, for every \(L>0\), the global source sequence \(g_L\)
of Theorem~\ref{thm:source-main} reconstructs every off-critical zero and multiplicity,
classifies every horizontal zero-free strip, and has exact false-RH negative-mass exponent
\(1+\delta_*\).
\end{theorem}

\subsection{Proof architecture}

Five mechanisms interact.
\begin{enumerate}[label=\arabic*.]
\item The fixed-width transform separates into two moving exponential edges.  After folding by
the functional equation, one edge is a Dirichlet coefficient
\(2\sum_{\Re w>0}A_L(w)^2n^w\); all remaining terms are uniformly bounded.
\item Its ordinary Dirichlet series is a sum of shifted zeta functions.  A
multiplicity-safe completion proves that every \(w\) with \(\Re w>0\) produces an uncancelled
simple pole at \(1+w\), with full-observable residue
\(2m(w)A_L(w)^2\).  The bounded remainder contributes a function holomorphic on
\(\Re s>1\), so pole positions and normalized residues reconstruct the off-critical divisor.
Landau's theorem for eventually nonnegative coefficients then converts any polynomial lower
budget into a horizontal zero bound.  Applying the same argument to an upper budget gives
two-sided oscillation.
\item A direct Guinand--Weil contour calculation evaluates the sharp window.  The
autocorrelation has compact support, so its prime side is finite.  Positivity and Abel summation
then identify negative Dirichlet abscissa with cumulative negative-mass growth.
\item A translated sector-to-disc map converts (1.18) into a Blaschke uniqueness theorem.
Lambert expansion and M\"obius inversion then recover every exponential moment, giving density
of the reflected derivative-Lambert family.
\item Differentiation of genuine pole-free local-factor changes identifies those Lambert
kernels as tangents.  Infinitesimal repetition upgrades linear density to a nonlinear orbit
law.  Odd-to-even normalization transports the orbit to the RH-equivalent Hankel cone.
\end{enumerate}

Sections \ref{sec:divisor}--\ref{sec:smooth} prove the global tomography statements, and
Section~\ref{sec:local} proves the local orbit and dichotomy.
Section \ref{sec:literature} gives the precise comparison with earlier localization and spectral
results.

\section{Centered divisor and two-edge transform}\label{sec:divisor}

The functional equation and real type of \(\Xi\) imply that \(\cZ\) is invariant under
\[
 w\mapsto-w,\qquad w\mapsto\overline w.
\tag{2.1}
\]
The classical zero count gives
\[
 \sum_{w\in\cZ}\frac1{|w|^2}<\infty.
\tag{2.2}
\]
There is no centered zero at \(w=0\), but defining \(A_L(0)=L\) makes all formulae entire.

\begin{lemma}[Nonvanishing and edge factorization]\label{lem:edge}
If \(L>0\) and \(\Re w\ne0\), then \(A_L(w)\ne0\).  For every \(w\in\C\),
\[
 A_L(-w)=e^{Lw}A_L(w)
\tag{2.3}
\]
and
\[
 F_{a,L}(w)^2=A_L(w)^2
 \left(e^{2aw}+2e^{Lw}+e^{-2(a-L)w}\right).
\tag{2.4}
\]
Uniformly on \(|\Re w|\le\frac12\),
\[
 |A_L(w)|\le\frac{1+e^{L/2}}{|w|}\qquad(w\ne0).
\tag{2.5}
\]
\end{lemma}

\begin{proof}
If \(A_L(w)=0\) and \(w\ne0\), then \(e^{-Lw}=1\).  Taking absolute values gives
\(e^{-L\Re w}=1\), hence \(\Re w=0\).  Identity (2.3) is immediate from (1.3), and expansion of
\(F_{a,L}^2\) gives (2.4).  Finally,
\(|1-e^{-Lw}|\le1+e^{-L\Re w}\le1+e^{L/2}\), proving (2.5).
\end{proof}

Equations (2.2) and (2.5) show that (1.5) converges absolutely and locally uniformly in \(a\).
The symmetries (2.1) show that its value is real.  On RH, \(w\) is purely imaginary and
\(F_{a,L}(w)\) is real after the usual Fourier identification.  Therefore every paired
contribution to \(Q_L(a)\) is a nonnegative square.

Let
\[
 \cZ_+=\{w\in\cZ:\Re w>0\},
 \qquad
 d_L(n)=2\sum_{w\in\cZ_+}A_L(w)^2n^w.
\tag{2.6}
\]
Multiplicity is retained.  Conjugation makes \(d_L(n)\) real.

\begin{proposition}[Bounded edge remainder]\label{prop:remainder}
For every \(L>0\) there is \(K_L<\infty\) such that
\[
 \left|q_L(n)-d_L(n)\right|\le K_L
\tag{2.7}
\]
at every \(L\)-admissible sample.
\end{proposition}

\begin{proof}
For \(\Re w>0\), retain \(A_L(w)^2e^{2aw}\) from (2.4).  For \(\Re w<0\), retain the last
term and fold it to the positive half-divisor by \(w\mapsto-w\).  At
\(a=\frac12\log n\), the retained sum is (2.6).  The middle term is independent of \(a\).
When \(a\ge L\), the discarded exponential on each open half-strip has modulus bounded by a
constant depending only on \(L\); on \(\Re w=0\), all exponential factors have bounded modulus.
By (2.5), the discarded contribution is bounded by
\[
 C_L\sum_{w\in\cZ}|w|^{-2},
\]
which is finite by (2.2).
\end{proof}

No rightmost zero is used in this decomposition.  The supremum \(\delta_*\) may fail to be
attained.

\section{Shifted-zeta carrier and exact pole divisor}\label{sec:carrier}

Assume \(\cZ_+\ne\varnothing\).  On \(\Re s>1+\delta_*\), absolute convergence gives
\[
 \sum_{n\ge1}\frac{d_L(n)}{n^s}
 =2C_L(s),
 \qquad
 C_L(s)=\sum_{w\in\cZ_+}A_L(w)^2\zeta(s-w).
\tag{3.1}
\]
Naive pole counting in (3.1) is insufficient when Xi zeros have multiplicity.  We therefore
record the cancellation-safe continuation.

\begin{proposition}[Multiplicity-safe pole tomography]\label{prop:poles}
The function \(C_L\) has a canonical meromorphic continuation to \(\Re s>1\).  Its poles are
exactly
\[
 s=1+w,\qquad w\in\cZ_+,
\tag{3.2}
\]
each pole is simple, and if \(w\) has analytic multiplicity \(m(w)\), then
\[
 \Res_{s=1+w}C_L(s)=m(w)A_L(w)^2\ne0.
\tag{3.3}
\]
\end{proposition}

\begin{proof}
The zero count also gives \(\sum_{w\in\cZ_+}|w|^{-3/2}<\infty\).  Multiply (3.1) by
\(\Xi(s-\frac12)\), fill the removable singularities term by term, and define
\[
 X_L(s)=\sum_{w\in\cZ_+}
 A_L(w)^2\,\Xi(s-\tfrac12)\zeta(s-w).
\tag{3.4}
\]
Standard vertical bounds for \(\zeta\) on closed substrips, (2.5), and the reciprocal
\(3/2\)-majorant give local normal convergence on \(\Re s>1\).  Hence \(X_L\) is holomorphic
there and
\[
 \widetilde C_L(s)=\frac{X_L(s)}{\Xi(s-\frac12)}
\tag{3.5}
\]
continues \(C_L\).

Fix a distinct \(w_0\in\cZ_+\) of multiplicity \(m\), and put \(p=1+w_0\).  Locally,
\[
 \Xi(s-\tfrac12)=(s-p)^m u(s),\qquad u(p)\ne0.
\tag{3.6}
\]
The \(m\) equal-support terms in (3.4) are
\[
 mA_L(w_0)^2(s-p)^{m-1}u(s)R(s-w_0),
\qquad
 R(z)=(z-1)\zeta(z),\quad R(1)=1.
\tag{3.7}
\]
Every term with support different from \(w_0\) is divisible by \((s-p)^m\).
Local normal convergence preserves that order floor.  Thus \(X_L\) has exact order \(m-1\) at
\(p\), while the denominator has order \(m\).  Division yields a simple pole with residue
\(mA_L(w_0)^2\).  Lemma \ref{lem:edge} makes the residue nonzero.  Away from (3.2), the
denominator of (3.5) is nonzero, so no other pole occurs.
\end{proof}

\begin{proposition}[Pole obstruction from an ordinary Dirichlet series]\label{prop:pole-obstruction}
Let \(b_n=d_L(n)+Cn^\alpha\) after deletion of a finite prefix, where \(\alpha\ge0\).
If \((b_n)\) is nonnegative and its ordinary Dirichlet series converges at a real point
\(\sigma\), then
\[
 \Re w\le\max\{\alpha,\sigma-1\}
\qquad(w\in\cZ_+).
\tag{3.8}
\]
The analogous assertion holds for \(Cn^\alpha-d_L(n)\).
\end{proposition}

\begin{proof}
On a far-right half-plane the series is
\[
 2C_L(s)+C\zeta(s-\alpha)-P(s),
\tag{3.9}
\]
with \(P\) a Dirichlet polynomial.  Its ordinary Dirichlet-series sum is holomorphic on
\(\Re s>\sigma\).  The identity theorem propagates agreement with (3.9) through every connected
pole-free component.  A shifted Xi pole \(1+w\) lying strictly to the right of both
\(\sigma\) and \(1+\alpha\) cannot be cancelled by the correction term.  Proposition
\ref{prop:poles} then contradicts holomorphy.  The proof for \(Cn^\alpha-d_L(n)\) is identical;
its shifted-Xi residues have the opposite sign but remain nonzero.
\end{proof}

\begin{proof}[Proof of Theorem \ref{thm:divisor-main}]
Set \(r_L(n)=q_L(n)-d_L(n)\).  Proposition \ref{prop:remainder} bounds \(r_L(n)\) on every
admissible tail; only finitely many indices are not admissible.  Therefore
\[
 R_L(s)=\sum_{n\ge1}\frac{r_L(n)}{n^s}
\]
converges absolutely and is holomorphic on \(\Re s>1\).  On
\(\Re s>1+\delta_*\),
\[
 \mathcal D_L(s)=2C_L(s)+R_L(s).
\]
Consequently
\[
 \mathcal D_L^\#(s)=2\widetilde C_L(s)+R_L(s)
\]
is the required continuation.  Proposition \ref{prop:poles} and holomorphy of \(R_L\) show that
its poles are exactly \(1+w\), all simple, with residue
\(2m(w)A_L(w)^2\).  Lemma \ref{lem:edge} makes the normalizing denominator nonzero, so the
residue quotient is precisely \(m(w)\).

The maps \(w\mapsto-w\) and \(w\mapsto\overline w\) recover the entire off-critical orbit from
its member with positive real part, including multiplicity.  If RH holds, \(\cZ_+\) is empty and
\(\mathcal D_L^\#=R_L\) is holomorphic on \(\Re s>1\).  Conversely, any off-critical zero has a
member in \(\cZ_+\), and the corresponding nonzero pole contradicts holomorphy.  This proves the
criterion and the divisor reconstruction.
\end{proof}

\section{Landau classification and two-sided oscillation}\label{sec:landau}

We use Landau's theorem in its classical form: an ordinary Dirichlet series with eventually
nonnegative real coefficients has a singularity at its finite abscissa of convergence.  Deleting
a finite prefix and adding \(Cn^\alpha\) preserve its applicability
\cite{Landau1905,Maurizi2010}.

\begin{lemma}[Scale-local Landau step]\label{lem:landau}
Suppose \(d_L(n)+Cn^\alpha\) is eventually nonnegative.  Then every
\(w\in\cZ_+\) satisfies \(\Re w\le\alpha\).  The same conclusion follows if
\(Cn^\alpha-d_L(n)\) is eventually nonnegative.
\end{lemma}

\begin{proof}
The series in (3.9) converges absolutely for \(\Re s>1+\delta_*\), so its abscissa
\(\sigma_c\) is at most \(1+\delta_*\).  If \(\delta_*>\alpha\), the carrier is holomorphic on a
disc about the real point \(1+\delta_*\): nontrivial Xi zeros have ordinate bounded away from
zero, and the local normal convergence used in Proposition \ref{prop:poles} is uniform there.
The correction \(\zeta(s-\alpha)\) is also holomorphic near that point.

Landau's theorem excludes \(\sigma_c=1+\delta_*\), so
\(\sigma_c<1+\delta_*\).  Choose
\[
 \max\{\sigma_c,1+\alpha\}<\sigma<1+\delta_*.
\]
The nonnegative series converges at \(\sigma\).  Proposition
\ref{prop:pole-obstruction} forces \(\Re w\le\sigma-1\) for every \(w\in\cZ_+\), contradicting
the definition of \(\delta_*\), since \(\sigma-1<\delta_*\).  Thus
\(\delta_*\le\alpha\).  The upper-budget case uses the second assertion of Proposition
\ref{prop:pole-obstruction}.
\end{proof}

\begin{proof}[Proof of Theorem \ref{thm:main}]
Under RH, the square interpretation following Lemma \ref{lem:edge} gives \(q_L(n)\ge0\).
The reciprocal-square majorant gives boundedness.  Thus RH implies all conditions in (1.7).

Conversely, suppose an admissible tail satisfies \(q_L(n)\ge B\).  Proposition
\ref{prop:remainder} gives \(d_L(n)\ge B-K_L\).  After adding a constant and deleting a finite
prefix, the carrier coefficients are nonnegative.  Lemma \ref{lem:landau} with \(\alpha=0\)
forces \(\cZ_+=\varnothing\).  Functional-equation symmetry then puts every zero on the critical
line.  This proves (1.7), and also shows directly that false RH makes \(q_L(n)\) cofinally
smaller than every real bound.

For the power upper bound, (2.5)--(2.6) and absolute summability give
\[
 |d_L(n)|\le
 2\left(\sum_{w\in\cZ_+}|A_L(w)|^2\right)n^{\delta_*}.
\tag{4.1}
\]
The bounded remainder is absorbed because \(n^{\delta_*}\ge1\).

If the negative inequality in (1.8) failed, then \(q_L(n)\ge-Cn^\alpha\) eventually.
Proposition \ref{prop:remainder} and Lemma \ref{lem:landau} would give
\(\delta_*\le\alpha\), contrary to \(\alpha<\delta_*\).  If the positive inequality failed,
then \(q_L(n)\le Cn^\alpha\) eventually; the upper-budget version of the same argument gives the
same contradiction.  This proves (1.8) and the exact threshold.

Finally, an eventual lower budget of exponent \(\alpha\) transfers to \(d_L\) and Lemma
\ref{lem:landau} gives \(\delta_*\le\alpha\).  Conversely, if all zeros satisfy
\(\Re\rho\le\frac12+\alpha\), then \(\delta_*\le\alpha\); (4.1) and the bounded remainder give
\(q_L(n)\ge-Cn^\alpha\).  Functional-equation symmetry accounts for the left half of the divisor.
This proves (1.9).
\end{proof}

\begin{remark}[No extremal zero premise]
The proof uses the support supremum only through convergence half-planes and inequalities.
Neither Proposition \ref{prop:poles} nor Lemma \ref{lem:landau} chooses a zero with
\(\Re w=\delta_*\).  The conclusion therefore remains valid if the transverse supremum is not
attained.
\end{remark}

\section{Sharp Guinand--Weil evaluation}\label{sec:explicit}

For the window (1.4), define its autocorrelation
\[
 R_{L,a}(t)=\int_{\R}h_{L,a}(x+t)h_{L,a}(x)\,dx.
\tag{5.1}
\]
It is an explicit continuous piecewise-linear function, vanishing for \(|t|>2a\).
Let
\[
 \widehat h_{L,a}(z)=F_{a,L}(iz)
\tag{5.2}
\]
and separate the completed factor
\[
 \cG(s)=\frac12s(s-1)\pi^{-s/2}\Gamma(s/2).
\tag{5.3}
\]
Put
\[
 \cP_{L,a}=
 \sum_{2\le k\le e^{2a}}
 \frac{\Lambda(k)}{\sqrt{k}}R_{L,a}(\log k)
\tag{5.4}
\]
and
\[
 \cA_{\infty,L,a}=
 \frac1\pi\int_{\R}
 \widehat h_{L,a}(x-i)^2
 \frac{\cG'}{\cG}\!\left(\frac32+ix\right)\,dx.
\tag{5.5}
\]
The sum in (5.4) is literal: it is over integers \(k\) satisfying the displayed bound.

\begin{theorem}[Sharp rectangular explicit formula]\label{thm:explicit}
If \(0\le L\le a\), the integral (5.5) converges absolutely and
\[
 Q_L(a)=-2\cP_{L,a}+\cA_{\infty,L,a}.
\tag{5.6}
\]
Both sides are real.
\end{theorem}

\begin{proof}
Let \(\Xi_C(z)=\Xi(\frac12+iz)\).  Integrate
\[
 \widehat h_{L,a}(z)^2\frac{\Xi_C'(z)}{\Xi_C(z)}
\tag{5.7}
\]
around rectangles with horizontal edges at imaginary parts \(\pm1\) and vertical edges
\(\Re z=\pm R\), choosing \(R\) away from zero coordinates.  The argument principle, with
analytic multiplicity, makes the residue sum tend to \(Q_L(a)\).  Evenness folds the boundary to
one horizontal and one vertical contribution.

We record the only estimate needed to remove the vertical edge.  If \(p_w=-iw\), the centered
Hadamard product gives
\[
 \frac{\Xi_C'(z)}{\Xi_C(z)}
 =\sum_{w\in\cZ}\frac{-z}{p_w(p_w-z)}.
\tag{5.8}
\]
The explicit quotient (1.3) gives
\(\widehat h_{L,a}(R+iy)=O_{L,a}(R^{-1})\) for \(|y|\le1\), after separating a possible nearby
divisor point by subtracting the analytic numerator at its ordinate.  Integrating the constant
part uses the logarithmic primitive of \((R+iy-p_w)^{-1}\).  Splitting at
\(|p_w|=3R\) yields
\[
 |V_R|\le C_{L,a}\left(
 R^{-1/2}\sum_{w\in\cZ}|p_w|^{-3/2}
 +R^{-1}\sum_{w\in\cZ}|p_w|^{-2}\right),
\tag{5.9}
\]
along a cofinal admissible sequence.  Both sums are finite, so \(V_R\to0\).

On the remaining edge,
\[
 \frac{\Xi_C'(x-i)}{\Xi_C(x-i)}
 =i\left(\frac{\cG'}{\cG}+\frac{\zeta'}{\zeta}\right)
   \left(\frac32+ix\right).
\tag{5.10}
\]
The Gamma--polar integrand is
\(O_{L,a}(\log(2+|x|)/(1+x^2))\), proving absolute convergence of (5.5).
The Euler series converges absolutely:
\[
 \frac{\zeta'}{\zeta}\!\left(\frac32+ix\right)
 =-\sum_{k\ge2}\frac{\Lambda(k)}{k^{3/2}}e^{-ix\log k}.
\tag{5.11}
\]
Fourier inversion for the compactly supported window gives
\[
 \frac1\pi\int_{\R}
 \widehat h_{L,a}(x-i)^2e^{-ixt}\,dx
 =2e^tR_{L,a}(t).
\tag{5.12}
\]
Termwise substitution of \(t=\log k\) in (5.11)--(5.12) gives
\[
 \frac1\pi\int_{\R}\widehat h_{L,a}(x-i)^2
 \frac{\zeta'}{\zeta}\!\left(\frac32+ix\right)\,dx
 =-2\sum_{k\ge2}\frac{\Lambda(k)}{\sqrt{k}}R_{L,a}(\log k).
\tag{5.13}
\]
Since \(R_{L,a}(t)=0\) for \(t>2a\), the last sum is (5.4).  Combining
(5.9)--(5.13) proves (5.6).  Conjugation followed by \(x\mapsto-x\) makes (5.5) real.
\end{proof}

\begin{remark}[Classical boundary]
The existence of an explicit formula for this compactly supported, piecewise \(C^1\)
autocorrelation also follows from the Guinand--Weil--Barner formula
\cite{Guinand1949,Barner1981}.  We do not claim a new general explicit formula.  The direct proof
fixes normalization, absolute convergence, and finite support in the conventions required by the
tomography theorem.  New content lies in coupling this evaluated sharp source to the fixed-width
logarithmic detector and its exact negative spectrum.
\end{remark}

At \(a=\frac12\log n\), define
\[
 g_L(n)=-2\cP_{L,\frac12\log n}
 +\Re\cA_{\infty,L,\frac12\log n}.
\tag{5.14}
\]
Theorem \ref{thm:explicit} gives (1.10).  In particular, the source coefficient is defined
without reference to the Xi divisor: it uses \(\Lambda\), the elementary window, and the
completed Gamma--polar factor only.

\begin{proof}[Proof of Corollary \ref{cor:source-divisor}]
For every \(n\ge N\), Theorem \ref{thm:explicit} gives \(g_L(n)=q_L(n)\).  Hence, on the initial
half-plane of convergence,
\[
 \mathcal G_{L,N}(s)
 =\mathcal D_L(s)-\sum_{1\le n<N}\frac{q_L(n)}{n^s}.
\]
The subtracted finite Dirichlet polynomial is entire.  Subtracting it from
\(\mathcal D_L^\#\) therefore gives a continuation with exactly the same poles, orders, and
residues.  The assertions now follow from Theorem \ref{thm:divisor-main}.
\end{proof}

\section{Negative-mass abscissa and cumulative law}\label{sec:mass}

For a nonnegative sequence \(c(n)\), define
\[
 \sigma_{\rm abs}(c)=
 \inf\left\{\sigma\in\R:
   \sum_{n\ge1}c(n)n^{-\sigma}<\infty\right\},
\tag{6.1}
\]
with value \(-\infty\) when \(c\) is eventually zero and the series is entire.

\begin{proposition}[Exact negative abscissa]\label{prop:abscissa}
For any admissible \(L,N\),
\[
 \sigma_{\rm abs}(g^-_{L,N})=
 \begin{cases}
 -\infty,&\mathrm{RH},\\
 1+\delta_*,&\neg\mathrm{RH}.
 \end{cases}
\tag{6.2}
\]
More generally, for every \(\alpha\ge0\),
\[
 \sigma_{\rm abs}(g^-_{L,N})\le1+\alpha
 \Longleftrightarrow
 \Xi(\rho)=0\Longrightarrow\Re\rho\le\tfrac12+\alpha.
\tag{6.3}
\]
\end{proposition}

\begin{proof}
Under RH, \(g_L(n)=q_L(n)\ge0\), so \(g^-_{L,N}\equiv0\).
Suppose RH is false and \(\sigma_{\rm abs}(g^-_{L,N})\le1+\alpha\).  Write
\(r_L(n)=q_L(n)-d_L(n)\); Proposition \ref{prop:remainder} makes \(r_L\) bounded, so its
Dirichlet series is holomorphic on \(\Re s>1\).  On the retained tail,
\[
 d_L(n)+r_L(n)+g^-_{L,N}(n)
 =q_L(n)+\max\{-q_L(n),0\}\ge0.
\tag{6.4}
\]
Fix \(\varepsilon>0\) and put \(\beta=\alpha+\varepsilon\).  The second correction in (6.4)
has an absolutely convergent Dirichlet series on \(\Re s>1+\beta\).  Let \(\sigma_c\) be the
abscissa of convergence of the nonnegative corrected series.  If
\(\sigma_c>1+\beta\), its meromorphic expression
\[
 2C_L(s)+\sum_{n\ge1}r_L(n)n^{-s}
       +\sum_{n\ge1}g^-_{L,N}(n)n^{-s}-P(s)
\]
is holomorphic near the real point \(\sigma_c\): both corrections are holomorphic there, \(P\)
is a Dirichlet polynomial, and the shifted-Xi poles of \(C_L\) are nonreal.  This contradicts
Landau's theorem.  Thus \(\sigma_c\le1+\beta\).  If some \(w\in\cZ_+\) satisfied
\(\Re w>\beta\), then \(1+w\) would lie in the open holomorphy half-plane of the nonnegative
Dirichlet series and both corrections, while Proposition \ref{prop:poles} gives \(C_L\) a
nonremovable pole there.  This is impossible.  Hence \(\Re w\le\alpha+\varepsilon\) for every
\(w\in\cZ_+\).  Letting \(\varepsilon\downarrow0\) gives \(\delta_*\le\alpha\).

Conversely, if \(\delta_*\le\alpha\), the power bound in Theorem \ref{thm:main} gives
\(g^-_{L,N}(n)=O(n^\alpha)\), so its series converges for \(\Re s>1+\alpha\).  This proves
the general equivalence.  Taking the infimum over \(\alpha\), and using \(\delta_*>0\) under
false RH, gives the dichotomy.
\end{proof}

\begin{lemma}[Positive Dirichlet series versus summatory growth]\label{lem:abel}
Let \(c(n)\ge0\), \(C(X)=\sum_{1\le n\le X}c(n)\), and \(\sigma\ge0\).  Then
\[
 \sigma_{\rm abs}(c)\le\sigma
 \Longleftrightarrow
 C(X)=O_\varepsilon(X^{\sigma+\varepsilon})
 \quad\text{for every }\varepsilon>0.
\tag{6.5}
\]
\end{lemma}

\begin{proof}
If the abscissa is at most \(\sigma\), fix \(\varepsilon>0\) and put
\(t=\sigma+\varepsilon/2\).  Positivity gives
\[
 C(X)\le X^t\sum_{n\ge1}c(n)n^{-t}
 =O(X^{\sigma+\varepsilon}).
\tag{6.6}
\]
Conversely, if the right side of (6.5) holds and \(s>\sigma\), use
\(\varepsilon=(s-\sigma)/2\).  Abel summation makes
\(\sum c(n)n^{-s}\) converge because the partial-sum exponent is strictly smaller than \(s\).
\end{proof}

\begin{proof}[Proof of Theorem \ref{thm:source-main}]
Identity (1.10) is Theorem \ref{thm:explicit} at logarithmic aperture.  Proposition
\ref{prop:abscissa} proves (1.11).  Equivalence of the half-plane bound and (1.12) is (6.3)
unpacked at every exponent \(1+\alpha+\varepsilon\).  Lemma \ref{lem:abel} gives (1.13).

Set \(\alpha=\delta_*\) to obtain the upper bound in (1.14).  If the smaller power
\(1+\delta_*-\eta\) were a summatory upper bound, Abel summation would give
\(\sigma_{\rm abs}(g^-_{L,N})\le1+\delta_*-\eta\), contradicting (1.11).  Taking
\(\alpha=0\) and using functional-equation symmetry gives (1.15).
\end{proof}

\begin{remark}[Meaning of the critical estimate]
Equation (1.15) is not advertised as an independently proved prime-number estimate.  The theorem
answers exactly what that estimate means: it is equivalent to RH.  Unconditional conclusions are
the source identity, all half-plane equivalences, exact false-RH exponent, and compulsory
two-sided oscillation.
\end{remark}

\section{Residual smooth two-edge profiles}\label{sec:smooth}

The sharp rectangle supplies the explicit arithmetic side.  We now show that fixed-width
completeness persists for a large set of smooth profiles.

Fix \(0<\eta<L-\eta\).  Let \(\mathcal S_{\eta,L}\) be the real Fr\'echet space of
\(C^\infty\) functions supported in \([\eta,L-\eta]\).  For \(f\in\mathcal S_{\eta,L}\), define
\[
 A_f(w)=\int_{\eta}^{L-\eta}f(x)e^{-wx}\,dx
\tag{7.1}
\]
and the translated two-edge transform
\[
 F_{a,f}(w)=e^{aw}A_f(w)+e^{-aw}A_f(-w).
\tag{7.2}
\]
Let \(Q_{f,L}(a)\) be the corresponding absolutely convergent Xi polarization.

\begin{lemma}[Generic simultaneous nonvanishing]\label{lem:generic}
For each \(w\in\cZ_+\), the map \(f\mapsto A_f(w)\) is a nonzero continuous real-linear map
into \(\C\).  Consequently
\[
 \cG_{\eta,L}=
 \bigcap_{w\in\cZ_+}\{f\in\mathcal S_{\eta,L}:A_f(w)\ne0\}
\tag{7.3}
\]
is dense and residual.
\end{lemma}

\begin{proof}
Continuity follows from the defining Fr\'echet seminorms.  The functional is nonzero: a smooth
bump supported in a sufficiently small subinterval has nonzero integral against the nowhere-zero
analytic weight \(e^{-wx}\).  Its kernel is therefore a proper closed real-linear subspace, so its
complement is open and dense.  The Xi divisor is countable.  The Baire category theorem applied
to the complete Fr\'echet space gives the conclusion.
\end{proof}

\begin{proposition}[Smooth carrier]\label{prop:smooth-carrier}
For every \(f\in\cG_{\eta,L}\), the fixed-width sequence \(Q_{f,L}(\frac12\log n)\) differs by a
bounded amount from
\[
 d_f(n)=2\sum_{w\in\cZ_+}A_f(w)^2n^w.
\tag{7.4}
\]
The Dirichlet series of \(d_f\) has a simple pole at each \(1+w\), with residue
\(2m(w)A_f(w)^2\ne0\).
\end{proposition}

\begin{proof}
The transform (7.2) has the same two exponential edges as (1.3).  Compact support away from the
endpoints gives rapid vertical decay of \(A_f\), so the discarded edges are absolutely summable
and uniformly bounded.  The shifted-zeta construction of Section \ref{sec:carrier} applies with
\(A_L(w)^2\) replaced by \(A_f(w)^2\).  Lemma \ref{lem:generic} prevents cancellation of every
shifted pole; the equal-support calculation (3.6)--(3.7) retains analytic multiplicity.
\end{proof}

\begin{proof}[Proof of Theorem \ref{thm:smooth-main}]
Under false RH, \(\cZ_+\ne\varnothing\).  For \(f\in\cG_{\eta,L}\), suppose the smooth samples
were eventually bounded below.  Proposition \ref{prop:smooth-carrier} transfers the lower bound
to \(d_f\).  Adding a constant and deleting a finite prefix gives an eventually nonnegative
ordinary Dirichlet series.  Landau's theorem and the noncancelled poles of Proposition
\ref{prop:smooth-carrier} give the same contradiction as Lemma \ref{lem:landau} with
\(\alpha=0\).  Hence the values are cofinally below every \(B\), proving (1.16).

Under RH, every smooth transform is square-summable on the centered divisor and its Weil
polarization is nonnegative.  Fix \(f\in\cG_{\eta,L}\).  If its samples were nonnegative or merely
eventually bounded below while RH was false, (1.16) with \(B=0\), or with the asserted lower
bound, would give a contradiction.  This proves the individual equivalences (1.17).  Applying
them to any member of the nonempty residual set gives the class-wide equivalences.
\end{proof}

The residual theorem answers the regularity concern completely: sharpness is needed for the
finite explicit source formula, not for fixed-width detector completeness.

\section{Universal local-factor orbit and RH information boundary}\label{sec:local}

We now prove the local half of the dichotomy.  Unlike Sections
\ref{sec:divisor}--\ref{sec:smooth}, this argument does not begin with the Xi divisor.

\subsection{Translated Blaschke uniqueness}

For an integer \(m\ge1\) and \(R\in\R\), let
\[
 \Sigma_{m,R}=
 \{R+w:w\ne0,\ |\arg w|<\pi/(2m)\}.
\tag{8.1}
\]
The conformal map
\[
 \psi_{m,R}(z)=
 R+\left(\frac{1+z}{1-z}\right)^{1/m}
\tag{8.2}
\]
maps the unit disc onto \(\Sigma_{m,R}\).  For \(\lambda>R+1\), the inverse image of
\(\lambda\) is
\[
 a_{\lambda,m,R}
 =\frac{(\lambda-R)^m-1}{(\lambda-R)^m+1}\in(0,1),
 \qquad
 1-|a_{\lambda,m,R}|
 =\frac2{(\lambda-R)^m+1}.
\tag{8.3}
\]

\begin{lemma}[Translated zero uniqueness]\label{lem:local-blaschke}
Assume (1.18) for one fixed \(m\) and all \(R\ge R_0\).  Suppose \(H\) is bounded and
holomorphic on every \(\Sigma_{m,R}\), \(R\ge R_0\), and
\[
 H(\lambda_j)=0\qquad(\lambda_j>R_0+1).
\tag{8.4}
\]
Then \(H(x)=0\) for every \(x\ge R_0+1\).
\end{lemma}

\begin{proof}
Fix \(x\ge R_0+1\) and take \(R=x-1\).  If \(H(x)\ne0\), then
\(H\circ\psi_{m,R}\) is a nonzero bounded holomorphic disc function.  Its distinct zeros include
all \(a_{\lambda_j,m,R}\) with \(\lambda_j>R+1\).  The zeros of such a function satisfy
\(\sum(1-|a_j|)<\infty\) \cite[Chapter II]{Garnett2007}.  Equations (1.18) and (8.3)
contradict this condition.  Hence \(H(x)=0\).
\end{proof}

The translation in (8.1) is what upgrades zero information at the spectral nodes to
vanishing on an entire positive ray.

\subsection{Lambert--M\"obius recovery}

Let \(K\) be compact and \(s:K\to\C\) continuous, with
\[
 \Re s(k)\ge\delta>0,\qquad |\Im s(k)|\le M.
\tag{8.5}
\]
For \(\lambda>0\), put
\[
 W_\lambda(k)=
 \lambda\frac{e^{-\lambda s(k)}}{(1-e^{-\lambda s(k)})^2}.
\tag{8.6}
\]

\begin{lemma}[Carrier completeness]\label{lem:carrier-complete}
Under (1.18), every coordinate monomial \(s^N\) belongs to the closed complex linear span of
the \(W_{\lambda_j}\) in \(C(K)\).
\end{lemma}

\begin{proof}
Let \(\ell\in C(K)^*\) annihilate every \(W_{\lambda_j}\), and define
\[
 F(z)=\ell(e^{-zs}),\qquad
 H(z)=\ell\!\left(\frac{e^{-zs}}{(1-e^{-zs})^2}\right).
\tag{8.7}
\]
The function \(F\) is entire.  Choose \(m\) so large that
\(\tan(\pi/(2m))<\delta/(M+1)\).  On every sufficiently right-translated sector
\(\Sigma_{m,R}\), condition (8.5) gives a positive uniform lower bound for
\(\Re(zs(k))\).  Thus \(H\) is bounded and holomorphic there.

Since \(\ell(W_{\lambda_j})=\lambda_jH(\lambda_j)=0\), Lemma
\ref{lem:local-blaschke} gives \(H(x)=0\) on a positive ray.  On that ray the
normally convergent Lambert expansion is
\[
 H(x)=\sum_{q\ge1}qF(qx).
\tag{8.8}
\]
Absolute convergence permits M\"obius inversion:
\[
\begin{aligned}
 \sum_{n\ge1}\mu(n)nH(nx)
 &=\sum_{n,q\ge1}\mu(n)nqF(nqx)\\
 &=\sum_{k\ge1}kF(kx)\sum_{n\mid k}\mu(n)
 =F(x).
\end{aligned}
\tag{8.9}
\]
Hence \(F\) vanishes on a ray and therefore identically.  Differentiation at zero gives
\[
 0=F^{(N)}(0)=(-1)^N\ell(s^N)
\tag{8.10}
\]
for every \(N\).  Hahn--Banach now proves the assertion.
\end{proof}

On \(K_r\), take \(s(k)=k\) and reflection \(\tau(s)=1-s\).  If a functional
annihilates every \(U_{\lambda_j}=W_{\lambda_j}-W_{\lambda_j}\circ\tau\), composing it with
\(f\mapsto f-f\circ\tau\) and using Lemma \ref{lem:carrier-complete} shows that it annihilates
every reflected polynomial
\[
 P(s)-P(1-s).
\tag{8.11}
\]
Taylor polynomials about \(1/2\), followed by oddification, approximate every holomorphic odd
germ uniformly on \(K_r\).  Finally Schwarz projection
\[
 f\longmapsto\frac{f(s)+\overline{f(\overline s)}}2
\tag{8.12}
\]
turns complex approximants into real linear combinations.  This proves the linear density in
Theorem \ref{thm:local-orbit}.

\subsection{From tangents to genuine local factors}

For \(q_\lambda(s)=e^{-\lambda s}\), the Banach-algebra map
\[
 a\longmapsto
 \Phi_{\lambda,a}
 =\lambda\frac{a q_\lambda}{1-aq_\lambda}
\tag{8.13}
\]
is holomorphic near \(a=1\), and
\[
 \left.\frac d{da}\Phi_{\lambda,a}\right|_{a=1}
 =\lambda\frac{q_\lambda}{(1-q_\lambda)^2}.
\tag{8.14}
\]
After reflection this gives
\[
 \left.\frac d{db}\Delta_{\lambda,b}\right|_{b=0}=U_\lambda.
\tag{8.15}
\]
The increments are genuine logarithmic derivatives.  If
\[
 V_{\lambda,b}(s)=
 \frac{1-(1+b)e^{-\lambda s}}{1-e^{-\lambda s}},
\tag{8.16}
\]
then on the pole-free locus (1.22),
\[
 \Delta_{\lambda,b}(s)
 =\frac d{ds}\log\!\bigl(V_{\lambda,b}(s)V_{\lambda,b}(1-s)\bigr).
\tag{8.17}
\]

Let \(q=\|q_\lambda\|_{K_r}<1\).  The inequality
\[
 |b|<\frac{1-q}{1+q}
\tag{8.18}
\]
implies \(\|(1+b)q_\lambda\|_{K_r}<1\), so every tangent has a nonempty
pole-free real parameter interval.  Differentiability at zero gives, for each \(c\in\R\),
\[
 (N+1)\Delta_{\lambda,c/(N+1)}
 \longrightarrow cU_\lambda
\tag{8.19}
\]
uniformly on \(K_r\).  The left side is an integer repetition of one permitted generator,
whose parameter eventually satisfies both (1.22) and any prescribed \(|b|<\epsilon\).
Thus the closure of the genuine increment group contains the real span of all
\(U_{\lambda_j}\), and hence the complete analytic fibre.

If \(A:C(K_r)\to Y\) is continuous, \(Y\) is Hausdorff, and \(A(x+g)=A(x)\) for every
allowed group element \(g\), then \(\{x:A(x)=A(0)\}\) is closed and contains the group.
It therefore contains its closure and the entire fibre.  This proves the one-point quotient.
Deleting finitely many frequencies preserves (1.18), proving all abstract assertions of
Theorem \ref{thm:local-orbit}.

\subsection{Prime logarithms}

Fix \(m,R\).  For all sufficiently large primes,
\[
 (\log p-R)^m+1\le2(\log p)^m\le2p.
\tag{8.20}
\]
Therefore
\[
 \sum_{\log p>R+1}\frac1{(\log p-R)^m+1}
 \ge\frac12\sum_{p\ {\rm sufficiently\ large}}\frac1p
 =\infty
\tag{8.21}
\]
by Euler's theorem.  Finite deletion changes neither divergence nor injectivity.  No prime
number theorem or zero-free region enters the local orbit result.

\subsection{The RH-equivalent moment sign}

The completed factorization and the chain rule give the continued prime--archimedean source
near the centre as
\[
 T_\zeta(z)=
 -i\left(
 -\frac{\zeta'}{\zeta}(\tfrac12+iz)
 -\frac{\cG'}{\cG}(\tfrac12+iz)\right)
 =\frac{\Xi_C'(z)}{\Xi_C(z)}.
\tag{8.22}
\]
It is odd and vanishes at zero.  Pairing the Hadamard product by the functional equation gives
\cite[Chapter 2]{Titchmarsh1986}
\[
\begin{aligned}
 G_\zeta(z)
 &=-\frac{T_\zeta(z)}z
 =\sum_\rho m_\rho\frac{\lambda_\rho}{1-z^2\lambda_\rho},\\
 \lambda_\rho&=-\frac1{(\rho-\frac12)^2},\qquad
 M_n=\sum_\rho m_\rho\lambda_\rho^{\,n+1}.
\end{aligned}
\tag{8.23}
\]
The inverse-square divisor sum is absolutely convergent.  Consequently, for
\(P(X)=\sum_jc_jX^j\in\R[X]\),
\[
 \sum_{i,j}c_ic_jM_{i+j}
 =\sum_\rho m_\rho\lambda_\rho P(\lambda_\rho)^2.
\tag{8.24}
\]

\begin{lemma}[Trace-class polynomial separation]\label{lem:trace-separation}
Let nonzero \(\lambda_j\in\C\) have positive integer weights \(w_j\) and
\(\sum_jw_j|\lambda_j|<\infty\).  If
\[
 \Re\sum_jw_j\lambda_jP(\lambda_j)^2\ge0
 \qquad(P\in\R[X]),
\tag{8.25}
\]
then every \(\lambda_j\) is real.
\end{lemma}

\begin{proof}
Suppose \(\alpha\) is nonreal and choose \(c\in\C\) with
\(\Re(\alpha c^2)<0\).  Only finitely many spectral values have modulus at least
\(|\alpha|\).  A product of real quadratic conjugate-pair factors gives a real polynomial
\(S\) which vanishes at all those values except \(\alpha,\overline\alpha\), and with
\(S(\alpha)\ne0\).

Put \(u=\alpha/|\alpha|\).  For each \(N\), real affine interpolation supplies a polynomial
\(L_N\) such that
\[
 L_N(\alpha)=\frac{c}{S(\alpha)u^N};
\tag{8.26}
\]
its coefficients stay bounded because \(|u|=1\).  Then
\[
 P_N(X)=S(X)L_N(X)(X/|\alpha|)^N
\tag{8.27}
\]
satisfies \(P_N(\alpha)=c\), vanishes at every unwanted outer value, is uniformly bounded on
\(|z|\le|\alpha|\), and tends to zero at every inner value.  Dominated convergence in
(8.25), using \(\sum_jw_j|\lambda_j|<\infty\), leaves a positive multiplicity sum of copies of
\(\Re(\alpha c^2)\), which is negative.  This contradiction proves the lemma.
\end{proof}

Under RH, \(\rho=\frac12+i\gamma\) and
\(\lambda_\rho=\gamma^{-2}>0\), so (8.24) is nonnegative.  Conversely, Lemma
\ref{lem:trace-separation} forces every \(\lambda_\rho\) to be real, while
\[
 \Im\lambda_\rho=
 \frac{2(\Re\rho-\frac12)\Im\rho}
 {|\rho-\frac12|^4}.
\tag{8.28}
\]
Nontrivial zeta zeros have nonzero imaginary part, hence every zero lies on the critical line.
Real-coefficient positivity is therefore equivalent to RH.  At each fixed degree the quadratic
form is continuous in the coefficient vector; rational approximation and clearing denominators
show that integer polynomials suffice.  This proves (1.24).

\subsection{Why the local quotient cannot carry the sign}

For any even Schwarz-real analytic germ \(g\), define its critical trace by
\[
 (\mathcal Tg)(s)=
 -(s-\tfrac12)g\!\left(-i(s-\tfrac12)\right).
\tag{8.29}
\]
It is functional-equation odd and Schwarz-real, and rotating back and dividing by the centred
coordinate recovers \(g\) exactly.  Theorem \ref{thm:local-orbit} therefore places
\(\mathcal Tg|_{K_r}\) in the closed pole-free prime-log orbit after every finite prime deletion.
Every continuous invariant reader \(A\) satisfies
\[
 A(\mathcal Tg|_{K_r})=A(0)
\tag{8.30}
\]
for every such \(g\).  But the zero germ satisfies all Hankel-square inequalities, whereas
\(g(z)=-1\) fails the test \(P=1\).  Thus no reader satisfying (8.30) can classify the sign,
proving the negative assertion of Theorem \ref{thm:dichotomy}.

The obstruction is topological rather than finite-jet.  Indeed, if
\(g(z)=\sum_{n\ge0}m_nz^{2n}\) is analytic on \(|z|<r_1\), then
\[
 g_N(z)=g(z)-(m_{2N}+1)z^{4N}
\tag{8.31}
\]
converges to \(g\) uniformly on every smaller disc, while its
\(z^{4N}\)-coefficient is \(-1\).  Hence the monomial test \(P=X^N\) fails for every
\(g_N\), and every corresponding trace still belongs to the same closed local-factor orbit.

The source sequence \(g_L(n)\) is not a function of this quotient.  It retains the indexed
weighted tail before local analytic continuation, and Corollary \ref{cor:source-divisor} proves
that this global presentation recovers all off-critical poles and multiplicities.  Theorem
\ref{thm:source-main} then recovers every horizontal support bound from its negative mass.
This proves the positive assertion and completes Theorem \ref{thm:dichotomy}.

\section{Relation to prior work and exact novelty boundary}\label{sec:literature}

The local theorem belongs first to the theory of complete analytic systems.  Classical
exponential completeness and M\"untz--Sz\'asz questions in holomorphic spaces have sharp
geometric and density criteria
\cite{Poltoratski2013,PelosoSalvatori2017,KhabibullinKudasheva2023}.  The family here is not a
pure exponential system: it is a reflected derivative-Lambert family, and its proof couples
translated sectorial Blaschke uniqueness to Lambert--M\"obius recovery.  The nonlinear theorem
then restricts the generators to logarithmic derivatives of invertible local factors.

This is also distinct from Voronin universality and generalized Euler-product universality.
Those theories use vertical translates of a global function and usually approximate
nonvanishing targets on compacta in a substrip
\cite{Voronin1975,LamzouriLesterRadziwill2018}.  Theorem \ref{thm:local-orbit} fixes the centre,
crosses the critical line, imposes functional-equation and Schwarz symmetry, permits vanishing
targets, and acts by finite local-factor modifications.  Neither result implies the other.

Hankel and moment reformulations of RH have substantial prior literature
\cite{Matiyasevich2007,Matiyasevich2017,BickelPascoeSargent2023}.  Equation (1.24) is used as an
exact sign target; novelty is not claimed for producing another equivalent criterion.  New
content is the incompatibility of that sign with the cofinite local-factor quotient, followed by
the explicit global sequence which survives and reconstructs the divisor.

Weil proved that positivity of the full explicit-formula quadratic functional is equivalent to RH
\cite{Weil1952}.  Guinand and Barner established broad explicit-formula frameworks that include
compactly supported autocorrelations of the present regularity
\cite{Guinand1949,Barner1981}.  Bombieri and Yoshida studied localized Hermitian forms and
fixed-support formulations \cite{Bombieri2000,Bombieri2003,Yoshida1992}.  In particular,
Bombieri's limiting negative-eigenvalue analysis assumes finitely many off-line zeros and permits
a limiting linear-dependence alternative.  Suzuki recast the Weil distribution through screw
functions and its Hilbert-space completion.  His 2026 localization theorem proves continuity of
the lowest localized eigenvalue and shows that false RH is equivalent to negativity for some
interval width \cite{Suzuki2023,Suzuki2025,Suzuki2026}.

Connes and Consani developed semilocal positivity and trace-formula structures; the CCM zeta
spectral triples use finite-prime self-adjoint models whose limiting identification remains a
separate issue \cite{Connes1999,ConnesConsani2021,ConnesConsani2023,CCM2025}.  Groskin gives an
exact finite Galerkin--Guinand--Weil dictionary and an archimedean tail certificate
\cite{Groskin2026}.  Those finite-rank results vary a coefficient vector or matrix size.  The
present theorem fixes one scalar profile before seeing the divisor and lets only its location
move.

The present contribution is narrower than those programs and orthogonal to their unresolved
operator limits:
\begin{enumerate}[label=\textup{(\arabic*)}]
\item every strongly sector-Blaschke-complete positive spectrum has a cofinite reflected-Lambert
completeness law and a dense pole-free finite-factor orbit;
\item for prime logarithms, the maximal continuous compact-open finite-factor-invariant quotient
of the analytic odd germ fibre is one point, so it cannot classify the RH-equivalent moment sign;
\item one arbitrarily fixed positive width and one logarithmically sampled translation geometry
has an ordinary Dirichlet continuation whose normalized polar divisor recovers every
off-critical Xi zero and its analytic multiplicity;
\item false RH forces cofinal two-sided oscillation at every subcritical transverse power;
\item the same sharp samples have an exact finite prime--Gamma realization, and every admissible
source tail has the identical polar divisor;
\item negative-part Dirichlet abscissa and cumulative mass recover every horizontal half-plane
bound exactly;
\item every individual profile in a dense residual smooth set is itself a complete detector and
has the same cofinal lower divergence.
\end{enumerate}

Items (1)--(2) identify an information boundary that does not occur in exponential or zeta
universality in this form.  Item (3) then supplies a global survivor rather than leaving the
no-go as methodological advice.  It strengthens the localization conclusion in a different
direction: failure is
witnessed not merely by some element of an expanding localized space, or by a minimizing
eigenvector, but cofinally by the translates of one predetermined rectangle of any positive
width, and it returns the complete off-critical divisor rather than only an extremal spectral
value.  Item (4) removes both the finite-off-line-zero condition and the
linear-dependence escape appearing in Bombieri's limiting spectral analysis.  Items (5)--(6)
then identify the same scalar witness at every sample with finite arithmetic data and recover the
complete horizontal support from its negative mass.

The general explicit formula, Weil positivity, Landau's theorem, Baire category, and the
positive-sequence Abel lemma are parent mechanisms, not novelty claims.  Nor are the Blaschke
condition, M\"obius inversion, an RH-equivalent Hankel criterion, or meromorphic divisor recovery
from a shifted-zeta series claimed as new mechanisms by themselves.  Priority is claimed for the
exact package: cofinite pole-free local-orbit collapse together with a single global
prime--Gamma sequence carrying the complete multiplicity-safe off-critical divisor and its exact
failure spectrum.

\section{Consequences and closed scope}\label{sec:scope}

The information question posed in the introduction has an exact two-level answer:
\begin{itemize}
\item Compact-open continuation to the centre has a trivial continuous quotient under
arbitrarily small pole-free finite-factor changes; this holds for every spectrum satisfying
(1.18), not only prime logarithms.
\item The RH-equivalent Hankel sign cannot descend to that quotient.  This is a complete no-go
for the simultaneous assumptions of local-germ dependence, continuity, and finite-factor
invariance, not a no-go for Hilbert--P\'olya, CCM, or trace formulae in general.
\item Every positive width works; no scale is selected from zero data.
\item One explicit sharp profile recovers every off-critical zero and multiplicity, detects RH,
detects every horizontal zero bound, and gives the exact size of every possible failure.
\item Every admissible sample equals a finite von Mangoldt sum plus an absolutely convergent
archimedean integral; any admissible source tail retains the complete off-critical polar divisor.
\item The sharp detector is not exceptional: individual completeness holds for every member of a
dense residual smooth class.
\end{itemize}

The global sequence answers the natural objection to a local no-go: it is an explicit surviving
object, not an unspecified stronger topology.  Its index \(n\), the finite von Mangoldt
incidence in each coefficient, and the coupled archimedean term retain precisely the weighted
tail presentation erased by the local quotient.

The result neither asserts RH nor declares an operator program impossible.  It proves that any
successful source mechanism must retain global weighted-tail information, react to finite local
changes, or be discontinuous in compact-open germ topology; the fixed-width prime--Gamma
sequence realizes the first alternative and is already an exact RH-complete detector.  Its
progress toward RH is the reduction from a large positivity class to one predetermined,
arithmetically evaluated scalar sequence, together with a complete law for the false case.  In
particular, the
critical estimate
\[
 \sum_{n\le X}g^-_{L,N}(n)=O_\varepsilon(X^{1+\varepsilon})
\]
is not assumed: Theorem \ref{thm:source-main} proves exactly that it is equivalent to RH.  Every
other assertion in the paper, including the source identity, pole noncancellation, cofinal
oscillation, and all exponent classifications, is unconditional.  No unresolved lemma or
conjectural input enters these conclusions.

\begin{thebibliography}{99}

\bibitem{Barner1981}
K.~Barner,
\emph{On A. Weil's explicit formula},
J. Reine Angew. Math. \textbf{323} (1981), 139--152.

\bibitem{BickelPascoeSargent2023}
K.~Bickel, J.~E.~Pascoe, and M.~Sargent,
\emph{Zero-free regions near a line},
Math. Z. \textbf{303} (2023), article 14.

\bibitem{Bombieri2000}
E.~Bombieri,
\emph{Remarks on Weil's quadratic functional in the theory of prime numbers, I},
Rend. Lincei Mat. Appl. \textbf{11} (2000), 183--233.

\bibitem{Bombieri2003}
E.~Bombieri,
\emph{A variational approach to the explicit formula},
Comm. Pure Appl. Math. \textbf{56} (2003), 1151--1164.

\bibitem{CCM2025}
A.~Connes, C.~Consani, and H.~Moscovici,
\emph{Zeta spectral triples},
arXiv:2511.22755, 2025.

\bibitem{Garnett2007}
J.~B.~Garnett,
\emph{Bounded Analytic Functions}, revised first edition,
Springer, New York, 2007.

\bibitem{Connes1999}
A.~Connes,
\emph{Trace formula in noncommutative geometry and the zeros of the Riemann zeta function},
Selecta Math. (N.S.) \textbf{5} (1999), 29--106.

\bibitem{ConnesConsani2021}
A.~Connes and C.~Consani,
\emph{Weil positivity and trace formula, the archimedean place},
Selecta Math. (N.S.) \textbf{27} (2021), article 77.

\bibitem{ConnesConsani2023}
A.~Connes and C.~Consani,
\emph{Spectral triples and zeta-cycles},
Enseign. Math. \textbf{69} (2023), 93--148.

\bibitem{Groskin2026}
A.~Groskin,
\emph{A finite Guinand--Weil dictionary and archimedean tail order for the truncated Weil
quadratic form},
arXiv:2607.02828, 2026.

\bibitem{Guinand1949}
A.~P. Guinand,
\emph{Fourier reciprocities and the Riemann zeta-function},
Proc. London Math. Soc. (2) \textbf{51} (1949), 401--414.

\bibitem{KhabibullinKudasheva2023}
B.~N. Khabibullin and E.~G. Kudasheva,
\emph{Completeness of exponential systems in spaces of functions on a compact or domain and
their area},
arXiv:2306.15899, 2023.

\bibitem{LamzouriLesterRadziwill2018}
Y.~Lamzouri, S.~Lester, and M.~Radziwi{\l}{\l},
\emph{An effective universality theorem for the Riemann zeta-function},
Comment. Math. Helv. \textbf{93} (2018), 709--765.

\bibitem{Landau1905}
E.~Landau,
\emph{\"Uber einen Satz von Tschebyschef},
Math. Ann. \textbf{61} (1905), 527--550.

\bibitem{Matiyasevich2007}
Y.~Matiyasevich,
\emph{Hidden life of Riemann's zeta function 2. Electrons and trains},
arXiv:0709.0028, 2007.

\bibitem{Matiyasevich2017}
Y.~V. Matiyasevich,
\emph{Riemann's hypothesis in terms of the eigenvalues of special Hankel matrices},
Proc. Steklov Inst. Math. \textbf{296} Suppl.~2 (2017), 78--91.

\bibitem{Maurizi2010}
B.~Maurizi,
\emph{Extending Landau's theorem on Dirichlet series with non-negative coefficients},
arXiv:1009.0228, 2010.

\bibitem{PelosoSalvatori2017}
M.~M. Peloso and M.~Salvatori,
\emph{Functions of exponential growth in a half-plane, sets of uniqueness and the
M\"untz--Sz\'asz problem for the Bergman space},
J. Geom. Anal. \textbf{27} (2017), 2570--2599.

\bibitem{Poltoratski2013}
A.~Poltoratski,
\emph{A problem on completeness of exponentials},
Ann. of Math. (2) \textbf{178} (2013), 983--1016.

\bibitem{Suzuki2023}
M.~Suzuki,
\emph{Aspects of the screw function corresponding to the Riemann zeta function},
arXiv:2206.03682, revised 2023.

\bibitem{Suzuki2025}
M.~Suzuki,
\emph{On the Hilbert space derived from the Weil distribution},
Canad. J. Math., published online 3 November 2025,
doi:10.4153/S0008414X25101739.

\bibitem{Suzuki2026}
M.~Suzuki,
\emph{Weil's quadratic form via the screw function},
arXiv:2606.09096, 2026.

\bibitem{Titchmarsh1986}
E.~C.~Titchmarsh,
\emph{The Theory of the Riemann Zeta-Function}, 2nd ed., revised by D.~R.~Heath-Brown,
Oxford University Press, 1986.

\bibitem{Weil1952}
A.~Weil,
\emph{Sur les ``formules explicites'' de la th\'eorie des nombres premiers},
Commun. S\'em. Math. Univ. Lund, volume dedicated to Marcel Riesz (1952), 252--265.

\bibitem{Voronin1975}
S.~M. Voronin,
\emph{Theorem on the universality of the Riemann zeta-function},
Math. USSR-Izv. \textbf{9} (1975), 443--453.

\bibitem{Yoshida1992}
H.~Yoshida,
\emph{On Hermitian forms attached to zeta functions},
in \emph{Zeta Functions in Geometry}, Adv. Stud. Pure Math. \textbf{21},
Mathematical Society of Japan, 1992, 281--325.

\end{thebibliography}

\end{document}
